% GENERATED FILE. Edit canonical lessons, then run npm run generate:books.
% canonical-source-hash: fnv1a64:8341453bbcfe66b2
% canonical-lessons: SA-C124-patah, SA-C124-sarovarah, SA-C124-naditiram, SA-C124-kokilah, SA-C124-catakah

\chapter{Cloth, Lake, Riverbank, Cuckoo, Sparrow}
\label{ch:sa-cloth-lake-riverbank-cuckoo-sparrow}
\hlchaptermodality{124}

\begin{chapteropening}
\textbf{By the end of this chapter:} \emph{I can say a cloth, a lake, a riverbank, a cuckoo and a sparrow.}

Complete the last lesson of chapter 124: I can say a cloth, a lake, a riverbank, a cuckoo and a sparrow.
\end{chapteropening}

\section[\texorpdfstring{\sk{पटः} (paṭaḥ)}{पटः (paṭaḥ)}]{\sk{पटः} (paṭaḥ) — a cloth}

\label{lesson:SA-C124-patah}



\begin{quote}
Before the new one: say the Sanskrit for a ship, then the Sanskrit for a train.
\end{quote}

\subsection*{You'll want to know: \sk{पटः}}
\textbf{\sk{पटः}} — \emph{paṭaḥ} — \textquotedblleft{}a cloth\textquotedblright{}.

\begin{cousinweb}
\emph{paṭa} \textquotedblleft{}woven cloth, a garment\textquotedblright{} is probably borrowed from Austroasiatic: CDIAL 7692 follows Przyluski in grouping it with \emph{karpāsa} \textquotedblleft{}cotton\textquotedblright{}, and thinks a Dravidian source less likely. Pali kept it as \emph{paṭa} \textquotedblleft{}cloth\textquotedblright{} (CDIAL 7692).
\end{cousinweb}

\begin{grammarlens}[title={What you've built}]
One more word for this chapter.
\end{grammarlens}

\subsection*{Your turn}
\begin{itemize}
\raggedright
  \item \emph{Say it:} \emph{paṭaḥ}
  \item \emph{Say it:} \emph{paṭaḥ}, once more
  \item \emph{Say it:} say \emph{paṭaḥ}
  \item \emph{Recall:} say the Sanskrit for a ship, then the Sanskrit for a train, then say \emph{paṭaḥ} again
\end{itemize}

\begin{tcolorbox}[breakable,colback=teal!4,colframe=teal!35!black,title={Before you move on}]
What does \sk{पटः} mean? (\textquotedblleft{}A cloth\textquotedblright{}.) Say it once more.
\end{tcolorbox}
\section[\texorpdfstring{\sk{सरोवरः} (sarovaraḥ)}{सरोवरः (sarovaraḥ)}]{\sk{सरोवरः} (sarovaraḥ) — a lake}

\label{lesson:SA-C124-sarovarah}



\begin{quote}
Before the new one: say the Sanskrit for a train, then the Sanskrit for a cloth.
\end{quote}

\subsection*{You'll want to know: \sk{सरोवरः}}
\textbf{\sk{सरोवरः}} — \emph{sarovaraḥ} — \textquotedblleft{}a lake\textquotedblright{}.

\begin{cousinweb}
\emph{sarovara} is built from \emph{saras} \textquotedblleft{}a lake, a pool\textquotedblright{}, from the root √sṛ \textquotedblleft{}flow\textquotedblright{} (CDIAL 13254), and \emph{vara} \textquotedblleft{}choice, best\textquotedblright{}: \textquotedblleft{}the finest of pools\textquotedblright{}, a lake deep enough for the lotus (MW).
\end{cousinweb}

\begin{grammarlens}[title={What you've built}]
One more word for this chapter.
\end{grammarlens}

\subsection*{Your turn}
\begin{itemize}
\raggedright
  \item \emph{Say it:} \emph{sarovaraḥ}
  \item \emph{Say it:} \emph{sarovaraḥ}, once more
  \item \emph{Say it:} say \emph{sarovaraḥ}
  \item \emph{Recall:} say the Sanskrit for a train, then the Sanskrit for a cloth, then say \emph{sarovaraḥ} again
\end{itemize}

\begin{tcolorbox}[breakable,colback=teal!4,colframe=teal!35!black,title={Before you move on}]
What does \sk{सरोवरः} mean? (\textquotedblleft{}A lake\textquotedblright{}.) Say it once more.
\end{tcolorbox}
\section[\texorpdfstring{\sk{नदीतीरम्} (nadītīram)}{नदीतीरम् (nadītīram)}]{\sk{नदीतीरम्} (nadītīram) — a riverbank}

\label{lesson:SA-C124-naditiram}



\begin{quote}
Before the new one: say the Sanskrit for a cloth, then the Sanskrit for a lake.
\end{quote}

\subsection*{You'll want to know: \sk{नदीतीरम्}}
\textbf{\sk{नदीतीरम्}} — \emph{nadītīram} — \textquotedblleft{}a riverbank\textquotedblright{}.

\begin{cousinweb}
\emph{nadītīra} is built from \emph{nadī} \textquotedblleft{}river\textquotedblright{}, which MW lists under the root √nad \textquotedblleft{}roar, sound\textquotedblright{}, and \emph{tīra} \textquotedblleft{}shore, bank\textquotedblright{} (CDIAL 5842): \textquotedblleft{}the bank of a river\textquotedblright{} (MW).
\end{cousinweb}

\begin{grammarlens}[title={What you've built}]
One more word for this chapter.
\end{grammarlens}

\subsection*{Your turn}
\begin{itemize}
\raggedright
  \item \emph{Say it:} \emph{nadītīram}
  \item \emph{Say it:} \emph{nadītīram}, once more
  \item \emph{Say it:} say \emph{nadītīram}
  \item \emph{Recall:} say the Sanskrit for a cloth, then the Sanskrit for a lake, then say \emph{nadītīram} again
\end{itemize}

\begin{tcolorbox}[breakable,colback=teal!4,colframe=teal!35!black,title={Before you move on}]
What does \sk{नदीतीरम्} mean? (\textquotedblleft{}A riverbank\textquotedblright{}.) Say it once more.
\end{tcolorbox}
\section[\texorpdfstring{\sk{कोकिलः} (kokilaḥ)}{कोकिलः (kokilaḥ)}]{\sk{कोकिलः} (kokilaḥ) — a cuckoo}

\label{lesson:SA-C124-kokilah}



\begin{quote}
Before the new one: say the Sanskrit for a lake, then the Sanskrit for a riverbank.
\end{quote}

\subsection*{You'll want to know: \sk{कोकिलः}}
\textbf{\sk{कोकिलः}} — \emph{kokilaḥ} — \textquotedblleft{}a cuckoo\textquotedblright{}.

\begin{cousinweb}
\emph{kokila} \textquotedblleft{}the Indian cuckoo, the koel\textquotedblright{} is imitative of the bird's call (MW), as Greek \emph{kokkyx} and English \emph{cuckoo} are (EtymDB). It lives on as Hindi \emph{koyal} (CDIAL 3483).
\end{cousinweb}

\begin{grammarlens}[title={What you've built}]
One more word for this chapter.
\end{grammarlens}

\subsection*{Your turn}
\begin{itemize}
\raggedright
  \item \emph{Say it:} \emph{kokilaḥ}
  \item \emph{Say it:} \emph{kokilaḥ}, once more
  \item \emph{Say it:} say \emph{kokilaḥ}
  \item \emph{Recall:} say the Sanskrit for a lake, then the Sanskrit for a riverbank, then say \emph{kokilaḥ} again
\end{itemize}

\begin{tcolorbox}[breakable,colback=teal!4,colframe=teal!35!black,title={Before you move on}]
What does \sk{कोकिलः} mean? (\textquotedblleft{}A cuckoo\textquotedblright{}.) Say it once more.
\end{tcolorbox}
\section[\texorpdfstring{\sk{चाटकः} (cāṭakaḥ)}{चाटकः (cāṭakaḥ)}]{\sk{चाटकः} (cāṭakaḥ) — a sparrow}

\label{lesson:SA-C124-catakah}



\begin{quote}
Before the new one: say the Sanskrit for a riverbank, then the Sanskrit for a cuckoo.
\end{quote}

\subsection*{You'll want to know: \sk{चाटकः}}
\textbf{\sk{चाटकः}} — \emph{cāṭakaḥ} — \textquotedblleft{}a sparrow\textquotedblright{}.

\begin{cousinweb}
MW and Apte record the sparrow as \emph{caṭaka}, with a short \emph{a}, and CDIAL 4571 calls it imitative of the bird's call, with similar words in Dravidian and Munda. It lives on as Hindi \emph{ciṛā} \textquotedblleft{}sparrow\textquotedblright{} and \emph{ciṛiyā} \textquotedblleft{}bird\textquotedblright{} (CDIAL 4571).
\end{cousinweb}

\begin{grammarlens}[title={What you've built}]
One more word for this chapter.
\end{grammarlens}

\subsection*{Your turn}
\begin{itemize}
\raggedright
  \item \emph{Say it:} \emph{cāṭakaḥ}
  \item \emph{Say it:} \emph{cāṭakaḥ}, once more
  \item \emph{Say it:} say \emph{cāṭakaḥ}
  \item \emph{Recall:} say the Sanskrit for a riverbank, then the Sanskrit for a cuckoo, then say \emph{cāṭakaḥ} again
\end{itemize}

\begin{tcolorbox}[breakable,colback=teal!4,colframe=teal!35!black,title={Before you move on}]
What does \sk{चाटकः} mean? (\textquotedblleft{}A sparrow\textquotedblright{}.) Say it once more.
\end{tcolorbox}
